% Source-aware proof for the repaired P06.1 v3 candidate.

\begin{proof}[Proof of
\cref{prop:p061-v3-single-shilov-factorization}]
Fix \(b\in V\). Extend scalars from \(K\) to \(\mathscr H(b)\), and
restrict the pullback of \(h\) to the analytic fibre \(W_b\). Scalar
extension and analytic pullback preserve classical pluriharmonicity by
Gubler--Rabinoff, Proposition 7.2(6),(7). Their affinoid
maximum--minimum principle, Proposition 7.3, says that both extrema occur
on the Shilov boundary of \(W_b\). Since that boundary is a singleton, the
pullback of \(h\) is constant on \(W_b\).

The underlying topological space of \(W_b\) is the set-theoretic fibre of
\(p\) over \(b\). Thus \(h\) is constant on every fibre of \(p\), and the
section identifies its value with \(h(s(b))\). This is precisely
\eqref{eq:p061-v3-single-shilov-factorization}.
\end{proof}

\begin{proof}[Proof of \cref{cor:p061-v3-single-shilov-rpl}]
By \cref{prop:p061-v3-single-shilov-factorization},
\(h|_W=p^*(s^*h)\). Gubler--Jell--Rabinoff, Lemma 5.4(1), gives pullback
stability of \(\mathbf R\)-PL functions for morphisms of good strictly
analytic spaces over a nontrivially valued field. All of those hypotheses
are now explicit, so the asserted regularity follows. Pullback also makes
\(s^*h\) classically pluriharmonic.
\end{proof}

\begin{proof}[Proof of
\cref{lem:p061-v3-contractive-standard-etale-tube}]
Completed tensor product commutes with these finite Tate presentations and
their closed ideals. With the norms specified in the statement, the quotient
presentation gives
\[
 D=B\langle t,u\rangle/
       \bigl(g_0(t),h_0(t)u-1\bigr).
\]
In particular,
\begin{equation}\label{eq:p061-v3-exact-central-relations}
 g_0(t)=0,
 \qquad h_0(t)u=1,
 \qquad \|t\|_D\leq1,
 \qquad \|u\|_D\leq1.
\end{equation}
Likewise \(\|T\|_{C_n},\|U\|_{C_n}\leq1\). The canonical maps
\(C\to C_n\) and \(D\to E_n\) are contractions for these quotient and
projective norms.

We first isolate a quantitative Banach estimate. Let \(R\) be a
non-Archimedean Banach algebra, let \(F\in R[X]\), and suppose
\(F'(z)\) is invertible. Write
\begin{equation}\label{eq:p061-v3-taylor-expansion}
 F(z+X)=F(z)+F'(z)X+\sum_{j=2}^{d}c_jX^j.
\end{equation}
If \(M\geq\|F'(z)^{-1}\|\), \(H\geq\max_{j\geq2}\|c_j\|\), and
\begin{equation}\label{eq:p061-v3-newton-inequalities}
                 MH r<1,
 \qquad \|F'(z)^{-1}F(z)\|\leq r,
\end{equation}
then
\[
 \Phi(\delta)=
 -F'(z)^{-1}F(z)
 -F'(z)^{-1}\sum_{j=2}^{d}c_j\delta^j
\]
is a strict contraction of the closed radius-\(r\) ball. Indeed, it maps
the ball to itself by the ultrametric inequality, and
\[
 \|\Phi(\delta)-\Phi(\delta')\|
 \leq MH r\,\|\delta-\delta'\|,
\]
because \(r<1\) and
\(\|\delta^j-(\delta')^j\|\leq
r^{j-1}\|\delta-\delta'\|\). Its unique fixed point gives the unique zero
of \(F\) in that ball. This also proves uniqueness directly: every zero in
the ball is a fixed point of \(\Phi\).

We now choose constants uniformly for all four root problems. The fixed
inverses
\[
                  g'(T)^{-1}\in C,
 \qquad g_0'(t)^{-1}\in D
\]
have bounded images in \(C_n\) and \(E_n\). Put
\[
 M=\max\bigl\{1,\|g'(T)^{-1}\|_C,
                    \|g_0'(t)^{-1}\|_D\bigr\}.
\]
The coefficient-contraction hypothesis and
\(\|t\|,\|T\|\leq1\) give
\begin{equation}\label{eq:p061-v3-polynomial-errors}
 \begin{aligned}
  \|g(t)\|_{E_n}&\longrightarrow0,&
  \|g'(t)-g_0'(t)\|_{E_n}&\longrightarrow0,\\
  \|g_0(T)\|_{C_n}&\longrightarrow0,&
  \|g_0'(T)-g'(T)\|_{C_n}&\longrightarrow0,\\
  \|h(t)-h_0(t)\|_{E_n}&\longrightarrow0,&
  \|h_0(T)-h(T)\|_{C_n}&\longrightarrow0.
 \end{aligned}
\end{equation}
After increasing \(n\), the two perturbed derivatives are invertible by
the Neumann series and all four derivative inverses have norm at most
\(M\).

Taylor expansion of \(g\) and \(g_0\) at \(t\) and at \(T\) produces
only finitely many coefficient expressions. Contractivity of every
structural map and the unit bounds on the two centres therefore give one
number \(H\geq1\) bounding every higher Taylor coefficient in all four
expansions, uniformly in \(n\); enlarge \(H\), if necessary, to cover the
fixed expansion of \(g_0\) at \(t\) in \(D\) as well. Similarly, there is
one \(L\geq1\) such
that both \(h\) and \(h_0\) are \(L\)-Lipschitz on each closed unit ball
appearing below.

Choose, once and for all,
\begin{equation}\label{eq:p061-v3-common-radius}
          0<r<1,
 \qquad M H r<1,
 \qquad Lr<1.
\end{equation}
Then increase \(n\) once more so that the two derivative perturbations
above are Neumann-small and
\begin{equation}\label{eq:p061-v3-common-error-bound}
 \begin{aligned}
 \|g'(t)^{-1}g(t)\|_{E_n}&\leq r,&
 \|g_0'(T)^{-1}g_0(T)\|_{C_n}&\leq r,\\
 \|h(t)-h_0(t)\|_{E_n}&<1,&
 \|h_0(T)-h(T)\|_{C_n}&<1.
 \end{aligned}
\end{equation}

Apply \eqref{eq:p061-v3-newton-inequalities} with this single \(r\) to
the following four triples:
\[
 (E_n,g,t),
 \qquad (E_n,g_0,t),
 \qquad (C_n,g_0,T),
 \qquad (C_n,g,T).
\]
For the first and third triples, \eqref{eq:p061-v3-polynomial-errors}
gives the required small Newton error and produces \(t_n\) and \(T_n^0\).
For the second and fourth triples the centre is already an exact root, so
the unique root is respectively \(t\) and \(T\). This proves, with an
actual common radius, all four uniqueness assertions in the statement.
The same estimate in \(D\) also makes \(t\) the unique zero of \(g_0\)
in its radius-\(r\) ball; this will be used after reduction.

The ultrametric inequality, the unit bounds on the centres, and \(r<1\)
give \(\|t_n\|_{E_n},\|T_n^0\|_{C_n}\leq1\). The localization factors have
the required bounds. For example,
\[
 \|h(t_n)-h_0(t)\|_{E_n}
 \leq\max\bigl\{L\|t_n-t\|,
                  \|h(t)-h_0(t)\|\bigr\}<1.
\]
Since \(u=h_0(t)^{-1}\) and \(\|u\|\leq1\), the element
\(h(t_n)=h_0(t)(1+u(h(t_n)-h_0(t)))\) is invertible and
\(\|h(t_n)^{-1}\|\leq1\). The same calculation with
\(U=h(T)^{-1}\) gives \(\|h_0(T_n^0)^{-1}\|\leq1\).

The universal property of the unit-radius Tate algebra now gives a
contractive \(A_n\)-homomorphism
\[
 \alpha_n:C_n\longrightarrow E_n,
 \qquad T\longmapsto t_n,
 \qquad U\longmapsto h(t_n)^{-1}.
\]
The displayed relations show that it factors through \(C_n\). Likewise,
the unit bounds define a contractive \(B\)-homomorphism
\[
 D\longrightarrow C_n,
 \qquad t\longmapsto T_n^0,
 \qquad u\longmapsto h_0(T_n^0)^{-1}.
\]
Together with \(A_n\to C_n\), the universal property of
\(A_n\widehat\otimes_BD\) gives the contractive
\(A_n\)-homomorphism \(\beta_n:E_n\to C_n\).

We use the two exact-root uniqueness statements, not merely the two
existence statements, to identify the composites. Since \(\alpha_n\) is
contractive,
\[
 \|\alpha_n(T_n^0)-t\|_{E_n}
 \leq\max\bigl\{\|T_n^0-T\|_{C_n},
                  \|t_n-t\|_{E_n}\bigr\}\leq r.
\]
The element \(\alpha_n(T_n^0)\) is a zero of \(g_0\), so assertion (2)
gives \(\alpha_n\beta_n(t)=t\). Similarly,
\[
 \|\beta_n(t_n)-T\|_{C_n}
 \leq\max\bigl\{\|t_n-t\|_{E_n},
                  \|T_n^0-T\|_{C_n}\bigr\}\leq r,
\]
and \(\beta_n(t_n)\) is a zero of \(g\); assertion (4) gives
\(\beta_n\alpha_n(T)=T\). The localization relations then fix \(u\) and
\(U\), and both composites already fix \(A_n\). These elements
topologically generate the two quotient Tate algebras, hence
\(\alpha_n\) and \(\beta_n\) are inverse.

Finally tensor the construction through \(A_n\to B\). The coefficient
errors vanish, and the image of \(t_n\) is a zero of \(g_0\) within
distance \(r\) of \(t\). The uniqueness estimate in \(D\) makes it
equal to \(t\); the localization inverse then equals \(u\). The same
argument applies to \(T_n^0\). Thus the reduced isomorphism is the
identity of \(D\).
\end{proof}

\begin{proof}[Proof of \cref{prop:p061-v3-etale-toric-product}]
Choose \(0<|\varpi|<1\). We first separate the zero-rank case. If
\(\operatorname{rank}(N)=0\), then
\[
 N=0,\qquad \rho=\{0\},\qquad S=\{0\},\qquad
 Z=\operatorname{Spec}K,
\]
and we put \(v=0\). If \(\operatorname{rank}(N)>0\), choose
\(v\in N\cap\operatorname{int}(\rho)\). In both cases put
\(\ell(m)=\langle m,v\rangle\). For \(k\geq1\), let
\begin{equation}\label{eq:p061-v3-weighted-toric-domain}
 U_k=\left\{z\in Z^{\mathrm{an}}:
  |\chi^m(z)|\leq|\varpi|^{k\ell(m)}
  \text{ for every }m\in S\right\}.
\end{equation}
If \(\operatorname{rank}(N)=0\), then every \(U_k\) is the one-point space
\(\{a\}=\mathcal M(K)\), and these domains are trivially cofinal among its
ordinary neighbourhoods. If \(\operatorname{rank}(N)>0\), then \(S\) is
sharp and finitely generated and \(\ell(m)>0\) for \(m\neq0\), so the
\(U_k\) are strict affinoid neighbourhoods of \(a\), cofinal among its
ordinary neighbourhoods. This is also the weighted-domain assertion in the
already production-verified
\cref{prop:p061-parameterized-toric-fixed-fibre}.

Let \(\xi\) be the scheme-theoretic support of \(x\). Fix an affine
target neighbourhood \(\operatorname{Spec}R\) of \(q(\xi)\). Stacks
Project, Tag 025C, gives, after shrinking around \(\xi\), an open immersion
from a source neighbourhood into
\[
 \operatorname{Spec}\bigl(R[\Theta,G'(\Theta)^{-1}]/(G(\Theta))\bigr),
\]
where \(G\in R[\Theta]\) is monic. Principal opens form a basis on this
affine standard-etale scheme. We may therefore choose a principal open
containing the image of \(\xi\) and contained in the image of the open
immersion. Represent its defining element by a polynomial
\(F\in R[\Theta]\), and put \(H=G'F\). The resulting actual source
neighbourhood is
\begin{equation}\label{eq:p061-v3-algebraic-standard-etale}
 \operatorname{Spec}
 \bigl(R[\Theta,H(\Theta)^{-1}]/(G(\Theta))\bigr),
\end{equation}
and \(G'(\Theta)\) is invertible there because the product \(G'F\) is
invertible. Thus \eqref{eq:p061-v3-algebraic-standard-etale} is genuinely
an open neighbourhood of \(\xi\) in the original source, not merely a
presentation of its local ring.

The production-verified product-neighbourhood lemma
\cref{lem:p061-product-open-tube} and cofinality of the \(U_k\) give a
strict affinoid neighbourhood \(V=\mathcal M(B)\) of \(y\) and an
integer \(m\) such that \(U_m\times_KV\) lies in the chosen affine target
open. Put
\[
 R_k=\mathcal O(U_k),
 \qquad A=R_m\widehat\otimes_KB,
 \qquad A_n=R_{m+n}\widehat\otimes_KB.
\]
We next record explicit norms rather than appeal to an unspecified
equivalent Banach norm. If \(\operatorname{rank}(N)=0\), take \(q=0\).
Then \(R_k=K\) for every \(k\), and hence
\[
                         A=A_n=B.
\]
Give all three algebras the fixed norm of \(B\); every structural map is the
identity and is a contraction. Suppose now that
\(\operatorname{rank}(N)>0\). Choose nonzero monoid generators
\(m_1,\ldots,m_q\) of \(S\), and set
\[
 X_i^{(k)}=\varpi^{-k\ell(m_i)}\chi^{m_i}.
\]
The algebra \(R_k\) is the quotient of
\(K\langle X_1^{(k)},\ldots,X_q^{(k)}\rangle\) by the binomial semigroup
relations; give it the quotient Gauss norm. This is exactly the semigroup
Gauss norm: choose once and for all one word in the generators representing
each \(s\in S\). Lifting a convergent semigroup series by those words does
not increase its Gauss norm, while the quotient map is contractive, so the
two inequalities between the quotient and semigroup norms are equalities.
Under restriction,
\begin{equation}\label{eq:p061-v3-normalized-restriction}
 X_i^{(m)}\longmapsto
 \varpi^{n\ell(m_i)}X_i^{(m+n)}.
\end{equation}
Thus \(R_m\to R_{m+n}\) is a contraction. Give \(A,A_n\) the quotient
projective tensor norms. The inclusion \(B\to A\), restriction
\(A\to A_n\), and evaluation at the fixed point
\(A,A_n\to B\), which sends every \(X_i^{(k)}\) to zero, are all
contractions. Together with the preceding zero-rank choice, this gives the
required compatible retractions in both cases.

We now construct the source affinoid and keep track of the chosen point.
Base change \(G,H\) from \(R\) to \(A\). Since the valuation is
nontrivial, choose \(c,d\in K^\times\) such that
\begin{equation}\label{eq:p061-v3-strict-source-bounds}
             |\Theta(x)|<|c|,
 \qquad |H(\Theta(x))|^{-1}<|d|.
\end{equation}
If \(e=\deg(G)\), put
\[
       T=c^{-1}\Theta,
 \qquad g(T)=c^{-e}G(cT),
 \qquad h(T)=dH(cT).
\]
Then \(g\) is monic. First form the closed analytic subspace cut out by
\(G(\Theta)=0\) in the relative closed disc \(|\Theta|\leq|c|\) over
\(\mathcal M(A)=U_m\times_KV\). Its algebra is
\(A\langle T\rangle/(g)\). Inside it
take the rational domain \(|1|\leq|dH(cT)|\). The two stages impose,
respectively,
\[
                 |\Theta|\leq|c|,
 \qquad |H(\Theta)|^{-1}\leq|d|.
\]
The usual rational-domain presentation shows, without any
unbounded algebraic localization, that their intersection is the strict
affinoid domain with algebra
\begin{equation}\label{eq:p061-v3-constructed-source-algebra}
 C=A\langle T,U\rangle/
       \bigl(g(T),h(T)U-1\bigr),
 \qquad U=(dH(\Theta))^{-1}.
\end{equation}
This domain maps into the analytification of
\eqref{eq:p061-v3-algebraic-standard-etale}, since the rational inequality
forces \(H\) to be invertible. The target affinoid
\(U_m\times_KV\) contains an ordinary open neighbourhood of \((a,y)\).
Moreover, the two inequalities in
\eqref{eq:p061-v3-strict-source-bounds} say precisely
\(|T(x)|<1\) and \(|U(x)|<1\). By continuity, the inverse image of the
target interior intersected with these two strict loci is an ordinary open
neighbourhood of \(x\) contained in \(\mathcal M(C)\). The derivative
\(g'(T)=c^{1-e}G'(cT)\) is invertible in \(C\). Equip \(C\) with the
displayed quotient Gauss norm.

Every base change remains concrete:
\begin{align}
 C_n=C\widehat\otimes_AA_n
   &=A_n\langle T,U\rangle/(g,hU-1),
 \label{eq:p061-v3-source-base-change}\\
 D=C\widehat\otimes_AB
   &=B\langle t,u\rangle/(g_0,h_0u-1).
 \label{eq:p061-v3-fixed-fibre-algebra}
\end{align}
The domains \(\mathcal M(C_n)\) are obtained by imposing the same two
bounds over \(U_{m+n}\times_KV\). The target point \((a,y)\) lies in the
ordinary interior of every one of these target affinoid neighbourhoods:
all nonconstant toric monomials vanish at \(a\), and \(V\) was chosen as a
neighbourhood of \(y\). The strict inequalities \(|T(x)|<1\) and
\(|U(x)|<1\) are unchanged. The preceding interior argument therefore
applies after every base change. Consequently every
\(\mathcal M(C_n)\) is an affinoid neighbourhood of \(x\), while
\(\mathcal M(D)\) is an affinoid neighbourhood of \(x\) in the fixed
fibre. This proves the compatibility that is not supplied merely by the
algebraic standard-etale theorem.

It remains to verify the coefficient contraction required by
\cref{lem:p061-v3-contractive-standard-etale-tube}. Let
\(J=\ker(A\to B)\). If \(\operatorname{rank}(N)=0\), then
\(A=A_n=B\), so \(J=0\). Thus every coefficient error
\(a^\circ=a-\iota\varepsilon(a)\) is identically zero in every \(A_n\), and
\eqref{eq:p061-v3-coefficient-contraction} holds with value zero for every
\(n\). This argument uses no maximum over the empty family of monoid
generators. We may therefore assume that \(\operatorname{rank}(N)>0\).
With the normalized presentation above, \(J\) is the
closed ideal generated by \(X_1^{(m)},\ldots,X_q^{(m)}\). Indeed, the
normalized tensor-product presentation identifies \(A\) with the semigroup
Tate algebra \(B\langle S\rangle\). In that algebra every
zero-constant-term series is obtained by
assigning each nonzero monomial to one chosen generator dividing it; the
resulting finitely many subseries converge in the Gauss norm. For each fixed
\(\delta\in J\), choose a representation
\(\delta=\sum_iX_i^{(m)}\delta_i\) with \(\delta_i\in A\). Then
\eqref{eq:p061-v3-normalized-restriction} and contractivity give
\begin{equation}\label{eq:p061-v3-explicit-kernel-decay}
 \|\delta\|_{A_n}
 \leq\max_i
   |\varpi|^{n\ell(m_i)}\|\delta_i\|_A
 \longrightarrow0.
\end{equation}
For every coefficient of \(g\) and \(h\), subtracting the lift through
\(B\to A\) of its image in \(B\) gives such a fixed element of \(J\).
There are only finitely many coefficients, so
\eqref{eq:p061-v3-coefficient-contraction} follows.
Consequently the coefficient-contraction hypothesis holds in both rank
cases.

Apply \cref{lem:p061-v3-contractive-standard-etale-tube}. For all large
\(n\), it gives
\[
 C_n\simeq A_n\widehat\otimes_BD
 \simeq R_{m+n}\widehat\otimes_KD.
\]
On affinoid spectra this is the claimed product
\(\Omega\simeq U_{m+n}\times_KD\), and reduction through the fixed-point
retraction identifies the fixed fibre with \(\{a\}\times D\). The
construction above shows that both sides are actual affinoid
neighbourhoods of \(x\), not only formal completions.

The base change of the standard etale source to \(V\) is quasi-smooth over
\(K\), because \(V\) is quasi-smooth and the algebraic chart is etale.
Quasi-smoothness is preserved on analytic domains, so its affinoid domain
\(D\) is quasi-smooth. Finally, for \(d_0\in D\), the fibre of
\(\Omega\to D\) has affinoid algebra obtained from the weighted semigroup
algebra of \(U_{m+n}\) by completed scalar extension to
\(\mathscr H(d_0)\). After the monomial normalization above it is the
semigroup Tate algebra
\[
 \mathscr H(d_0)\langle S\rangle
 =\left\{\sum_{s\in S}a_sX^s:
               |a_s|\longrightarrow0\text{ off finite subsets}\right\}
\]
with norm \(\max_s|a_s|\). This norm is multiplicative: normalize two
nonzero series by coefficients of maximal norm and reduce their product to
the domain \(\widetilde{\mathscr H(d_0)}[S]\); the latter is a domain because
\(S\) is cancellative and torsion free. The multiplicative Gauss norm is a
Berkovich point attaining the supremum of every analytic function. Its
singleton is therefore a boundary, and the minimal nonempty Shilov boundary
is exactly that singleton.
\end{proof}

\begin{proof}[Proof of \cref{thm:p061-v3-etale-toroidal-ph-rpl}]
Let \(\mathcal U\) be the ordinary open domain of the given function.
First use \cref{prop:p061-v3-etale-toric-product} to obtain a product
neighbourhood \(U\times_KD\) of \(x\). Inside its intersection with
\(\mathcal U\), the product-neighbourhood lemma
\cref{lem:p061-product-open-tube}, followed by a deeper weighted toric
domain, gives a product \(U'\times_KD'\) still containing an ordinary
neighbourhood of \(x\), with \(D'\) a strict affinoid domain in \(D\).
Thus \(D'\) remains good, strict, and quasi-smooth.

The projection \(p:U'\times_KD'\to D'\) has a fixed section and
singleton-Shilov analytic fibres. By
\cref{prop:p061-v3-single-shilov-factorization}, the function factors as
\[
                         h=p^*g,
 \qquad g\in\operatorname{PH}(D').
\]
The exact production-verified quasi-smooth dependency
\cref{thm:p061-quasismooth-ph-rpl} makes \(g\) locally
\(\mathbf R\)-PL and Gubler--Jell--Rabinoff harmonic. The hypotheses of
\cref{cor:p061-v3-single-shilov-rpl} are satisfied, so \(h\) is locally
\(\mathbf R\)-PL. Since both \(h\) and \(-h\) are classically psh,
Gubler--Rabinoff, Theorem 4.3, makes both signs semipositive; this is
Gubler--Jell--Rabinoff harmonicity.
\end{proof}

\begin{proof}[Proof of
\cref{cor:p061-v3-zariski-etale-toric-charts}]
Fix an analytic point and a chart. If its image in the split toric factor
lies in the orbit of a face \(\tau\), the exact monoid localization in the
production-verified proof of
\cref{thm:p061-split-toric-all-orbits} gives a principal neighbourhood
isomorphic to
\[
                         Z_\tau\times_KT,
\]
where \(Z_\tau\) is the affine toric variety of a full-dimensional pointed
cone, viewed at its fixed point, and \(T\) is a split torus. The factor
\(T\) times the original smooth factor is smooth. At the dense torus orbit,
the pointed factor has rank zero, so \(Z_\tau=\operatorname{Spec}K\); this is
exactly the zero-rank branch proved above. Around the chosen source point,
replace the restricted ordinary etale chart by the affine
standard-etale neighbourhood supplied by Stacks Project, Tag 025C; in
particular it is separated and of finite type. Apply
\cref{thm:p061-v3-etale-toroidal-ph-rpl}. The
conclusion is local on the source, and the Zariski chart domains cover
\(\mathcal X^{\mathrm{an}}\).
\end{proof}
